\section{The Fire That Is Not Extinguished}

The loss is the chief pain, but it is not the only one, and the Lord did not leave his disciples to imagine that it was. In the ninth chapter of Mark he spoke of the hand, the foot, and the eye that cause a man to fall, and three times he closed the warning with the same words:

\begin{quote}
And if thy hand scandalize thee, cut it off: it is better for thee to enter into life, maimed, than having two hands to go into hell, into unquenchable fire: Where their worm dieth not, and the fire is not extinguished. And if thy foot scandalize thee, cut it off: it is better for thee to enter lame into life everlasting than having two feet to be cast into the hell of unquenchable fire: Where their worm dieth not, and the fire is not extinguished. And if thy eye scandalize thee, pluck it out: it is better for thee with one eye to enter into the kingdom of God than having two eyes to be cast into the hell of fire: Where their worm dieth not, and the fire is not extinguished.\endnote{Mark 9:42--47, Douay--Rheims (Challoner); the refrain stands at verses 43, 45 and 47.}
\end{quote}

Augustine heard the repetition as a cry. ``He did not shrink from using the same words three times over in one passage. And who is not terrified by this repetition, and by the threat of that punishment uttered so vehemently by the lips of the Lord Himself?''\endnote{Augustine, \emph{The City of God} XXI.9, trans.\ Marcus Dods (Edinburgh: T. \& T. Clark, 1871).} The words were taken up by the Church and laid on her lips. The Creed names ``everlasting fire''; the \emph{Catechism}, quoting the councils, names ``eternal fire''; and Paul VI ended his profession of faith with ``the fire that is not extinguished.''

\subsection{A real fire}

Augustine allowed a choice about one thing in Christ's refrain, and it was not the fire. Some, he reports, refer both the fire and the worm to the spirit; others take the fire as bodily and the worm as the gnawing grief of the soul. He himself found it ``easier to understand both as referring to the body than to suppose that neither does,'' and he let his reader choose concerning the worm: ``Let each man, then, choose which he will, whether he thinks that the worm is real and pertains to the body, or that spiritual things are meant by bodily representations, and that it belongs to the soul.'' But of the fire and the bodies of the damned he left no choice at all: ``only, this we believe about those future bodies, that they shall be such as shall certainly be pained by the fire.''\endnote{Augustine, \emph{City of God} XXI.9.} In the next chapter he goes further. The fire prepared for the devil and his angels is one fire, and it will torment both men and demons, the bodies of the one and the spirits of the other: ``But that hell, which also is called a lake of fire and brimstone, will be material fire, and will torment the bodies of the damned, whether men or devils\ldots{} One fire certainly shall be the lot of both, for thus the truth has declared.''\endnote{Ibid., XXI.10.}

Saint Gregory the Great taught the same in the fourth book of his \emph{Dialogues}, which he gave to the life of the soul after death. When his deacon Peter wonders how a bodily fire can hold a soul that has no body, Gregory answers that a spirit is held in the living body which it quickens, and so may be held in the fire in which it continually dies; and he adds that the soul suffers by seeing the fire and also by feeling it, ``for the soul by seeing of the fire is afflicted, and burned it is, in that it seeth itself to be burned.'' He appeals to the rich man of the parable, who begged for a drop of water ``because I am tormented in this flame.''\endnote{Gregory the Great, \emph{Dialogues} IV.29, trans.\ P.W. (1608), re-edited by Edmund G. Gardner (London: Philip Lee Warner, 1911), tertullian.org, accessed 16 October 2026; cf.\ Luke 16:24.} And when Peter asks whether there is one fire or many, Gregory gives the answer that every later teacher repeats:

\begin{quote}
The fire of hell is but one: yet doth it not in one manner torment all sinners. For every one there, according to the quantity of his sin, hath the measure of his pain. For as, in this world, many live under one and the same sun, and yet do not alike feel the heat thereof: for some be burnt more, and some less: so in that one fire, divers manners of burning be found, for that which in this world diversity of bodies doth, that in the next doth diversity of sins: so that although the fire be there all alike, yet doth it not in one manner and alike burn and torment them that be damned.\endnote{Gregory, \emph{Dialogues} IV.43 (chapter numbering of the 1911 edition). Gregory gives the same image of the one sun in \emph{Morals on the Book of Job} IX.65.98.}
\end{quote}

\noindent This is the inequality of punishments that the Second Council of Lyons professed, seen from within: one fire, burning each according to what he has done.

The schoolmen carried the Fathers' teaching forward without loosening it. The Supplement to the \emph{Summa}, reviewing the opinions about the fire, concludes that ``whatever we may say of the fire that torments the separated souls, we must admit that the fire which will torment the bodies of the damned after the resurrection is corporeal, since one cannot fittingly apply a punishment to a body unless that punishment itself be bodily.'' It holds, moreover, ``that the fire of hell is of the same species as the fire we have, so far as the nature of fire is concerned,'' while having ``certain properties differing from our fire, for instance that it needs no kindling, nor is kept alive by fuel.''\endnote{\emph{Summa theologiae}, Supplement, q.~97, aa.~5--6, trans.\ Fathers of the English Dominican Province (London: Burns Oates and Washbourne, 1922).} Saint John Damascene, as the Supplement cites him, had written that the ungodly will be cast into everlasting fire, ``not material fire, such as that which we have, but such as God knoweth''; the Supplement answers that he ``does not absolutely deny that this fire is material, but that it is material as our fire, since it differs from ours in some of its properties.''\endnote{Supplement, q.~97, a.~5, objection 1 and reply, citing John Damascene, \emph{On the Orthodox Faith} IV.} Long before, the Roman apologist Minucius Felix had described that difference: ``Nor is there either measure or termination to these torments. There the intelligent fire burns the limbs and restores them, feeds on them and nourishes them.''\endnote{Minucius Felix, \emph{Octavius} 35, trans.\ R. E. Wallis, \emph{The Ante-Nicene Fathers}, vol.~4, New Advent, accessed 16 October 2026.}

The Fathers and Doctors commonly hold, then, that it is a true fire, acting on the soul by the power of God and on the risen body as fire acts on a body, and unlike ours in that it burns without consuming and needs no fuel. What none of them holds is that the word is empty. Christ chose it three times in one breath, and he chose it to be feared.

\subsection{The meditation of the fifth exercise}

The saints did not leave the fire to the theologians. They placed it before the eyes of the soul, deliberately and in order, as the Church's own way of praying. The classic form is the fifth exercise of the first week of the \emph{Spiritual Exercises} of Saint Ignatius, which its author calls simply ``a meditation on hell.'' It is short, and it should be read whole:

\begin{quote}
\textsc{First Prelude.} The first Prelude is the composition, which is here to see with the sight of the imagination the length, breadth and depth of Hell.

\textsc{Second Prelude.} The second, to ask for what I want: it will be here to ask for interior sense of the pain which the damned suffer, in order that, if, through my faults, I should forget the love of the Eternal Lord, at least the fear of the pains may help me not to come into sin.

\textsc{First Point.} The first Point will be to see with the sight of the imagination the great fires, and the souls as in bodies of fire.

\textsc{Second Point.} The second, to hear with the ears wailings, howlings, cries, blasphemies against Christ our Lord and against all His Saints.

\textsc{Third Point.} The third, to smell with the smell smoke, sulphur, dregs and putrid things.

\textsc{Fourth Point.} The fourth, to taste with the taste bitter things, like tears, sadness and the worm of conscience.

\textsc{Fifth Point.} The fifth, to touch with the touch; that is to say, how the fires touch and burn the souls.\endnote{Ignatius of Loyola, \emph{Spiritual Exercises}, First Week, Fifth Exercise, trans.\ Elder Mullan (New York: P. J. Kenedy \& Sons, 1914), pp.~45--46.}
\end{quote}

The senses that sinned are the senses that are made to consider, one by one, what sin earns. Ignatius asks first for a grace, and he names the grace exactly: that if love should fail, fear at least may keep the soul from sin. The meditation is the servant of that petition. And it does not end in the fire. It ends in a colloquy with Christ, in which the soul recalls those who are lost, ``some because they did not believe the Coming, others because, believing, they did not act according to His Commandments,'' and then turns to give thanks: ``with this I will give Him thanks that He has not let me fall into any of these divisions, ending my life. Likewise, I will consider how up to now He has always had so great pity and mercy on me. I will end with an \textsc{Our Father}.''\endnote{Ibid., Fifth Exercise, Colloquy, p.~46.}

Francis de Sales gives the same meditation to beginners in the devout life, and he composes the place in a single image: ``Picture to yourself a dark city, reeking with the flames of sulphur and brimstone, inhabited by citizens who cannot get forth.'' The lost, he says, suffer ``in every sense and every member; and that because having used their members and senses for sin, it is just that through them they should suffer now.''\endnote{Francis de Sales, \emph{Introduction to the Devout Life} I.15, Preparation and first Consideration, p.~41.} Alphonsus walks through the same senses and comes at last to the touch, the sense that the fire torments most: ``They shall have an abyss of fire below, an abyss of fire above, and an abyss of fire on every side.''\endnote{Alphonsus, \emph{Preparation for Death}, Consideration XXVI, second point, p.~319.}

Such meditations are best made slowly, and best ended as Ignatius ends them, with thanks. The man who has looked at the fire honestly knows that he was not kept from it by his own strength. He was kept by the pity and mercy of Christ, ``up to now,'' and he will be kept tomorrow by the same.
